\documentclass[conference]{IEEEtran}

\usepackage[utf8]{inputenc}
\usepackage{amsmath,amssymb,amsfonts,bm}
\usepackage{booktabs}
\usepackage{multirow}
\usepackage{array}
\usepackage{enumitem}
\usepackage{algorithm}
\usepackage{algorithmic}
\usepackage{graphicx}
\usepackage{xcolor}
\usepackage{url}
\usepackage[hidelinks]{hyperref}

\DeclareMathOperator*{\argmin}{arg\,min}
\DeclareMathOperator*{\argmax}{arg\,max}
\DeclareMathOperator{\sigmoid}{sigmoid}
\DeclareMathOperator{\LCVaR}{LCVaR}
\DeclareMathOperator{\TopK}{TopK}

\newcommand{\E}{\mathbb{E}}
\newcommand{\cA}{\mathcal{A}}
\newcommand{\cG}{\mathcal{G}}
\newcommand{\cP}{\mathcal{P}}
\newcommand{\cR}{\mathcal{R}}
\newcommand{\cZ}{\mathcal{Z}}
\newcommand{\rvn}{\mathrm{OC\mbox{-}MERO}}
\newcommand{\sr}{\mathrm{SRCA}}
\newcommand{\cnro}{\mathrm{CNRO}}
\newtheorem{proposition}{Proposition}

\title{OC-RAP: Observation-Consistent Recovery Planning via Support--Reserve Complementarity}
\author{Anonymous Authors}

\begin{document}
\maketitle

\begin{abstract}
A safe autonomous-driving planner must preserve nominal utility in ordinary traffic, avoid entering unrecoverable states near contact, and continue improving safety after contact. The central difficulty is informational: branch-wise planning may certify a candidate by using different recovery maneuvers for latent futures that remain indistinguishable to the deployed vehicle. We formulate this mismatch as an \emph{oracle-to-deployable recoverability gap} and introduce \textbf{OC-RAP}, a single regime-agnostic recovery planner. OC-RAP pairs candidate prefixes with actuator-realizable recovery continuations and uses \textbf{OC-MERO} to evaluate each fixed recovery option over observation-compatible latent roots before nested lower-tail aggregation. A non-compensatory support--reserve admission then distinguishes support establishment from signed reserve/debt, while \textbf{RIFA} keeps absolute feasibility isolated from relative ranking. Our current mechanism analysis further shows that recovery orientation cannot be treated as a candidate-invariant scene direction: different candidates can activate different limiting interactions. We therefore introduce a constraint-native, nominal-zero audit based on the signed candidate-induced change of active clearance, with dimension-matched nearest-agent controls and candidate-identity shuffling. This yields one deployable information pattern and one signed recovery semantics across Safe, Near-Contact, and Contact, without regime routers, privileged future input, or regime-specific thresholds.
\end{abstract}

\section{Introduction}
\label{sec:introduction}

Collision avoidance alone is an incomplete notion of closed-loop driving safety. A planner must preserve useful fallback actions before contact and continue making safety-improving decisions after contact, where stabilization, re-contact avoidance, secondary-collision prevention, and safe re-entry remain relevant~\cite{koopman2024anatomy,wang2023integrated,parseh2023motion,wang2024post,yang2025post,ghosh2026post}. We therefore treat \emph{recoverability} as an action-level planning primitive rather than a terminal collision label.

This view exposes an information-pattern failure. Suppose multiple latent futures remain observationally indistinguishable after a candidate prefix. A branch-wise oracle may brake in one future and laterally escape in another, yet the deployed vehicle cannot condition on the hidden future identity. Existence of a branch-specific recovery is therefore weaker than existence of a \emph{deployable} recovery. OC-RAP addresses this oracle-to-deployable gap by requiring recovery decisions to be measurable with respect to the post-prefix observation.

A second challenge is that recoverability is signed and state-dependent. Before a common executable recovery exists, a useful action must establish support. Once support exists, the action must preserve that support while increasing positive reserve or repaying negative recovery debt. Near-Contact and Contact should therefore not be implemented as separate experts: they are different signs/states of the same recoverability object. This motivates a single policy with shared decision boundaries rather than a regime router.

A third challenge is causal representation. Our recent audits show that neither a fixed latent axis nor a candidate-invariant scene-conditioned direction transfers consistently across Near/Contact and development/certificate populations. A candidate changes which interaction is limiting. Consequently, the invariant object should be the \emph{candidate-conditioned constraint-response rule}, not a scene-wide orientation. The current V48.111 audit operationalizes this principle through a nominal-zero signed clearance response, $\Delta h_i^a(t)=h_i^a(t)-h_i^{a_0}(t)$, and compares candidate-specific active-agent geometry with an equal-dimensional nearest-agent control. This audit does not enter the runtime policy unless its preregistered cross-population gates are satisfied; it is an identifiability-before-capacity mechanism for deciding what representation may be promoted.

Finally, relative evidence that one candidate improves over nominal is not equivalent to absolute evidence that it is deployably recoverable. We therefore use \textbf{Role-Isolated Feasibility Admission (RIFA)}: the frozen proposal/ranking path is separated from the absolute recoverability predicate, preventing a large relative score from compensating for absolute infeasibility.

Our contributions are fourfold:
\begin{enumerate}[leftmargin=*]
    \item \textbf{Observation-consistent deployable recoverability.} We formalize the oracle-to-deployable recoverability gap and introduce OC-MERO, which evaluates a fixed recovery option over observation-compatible latent roots before any option selection and then performs nested lower-tail aggregation.
    \item \textbf{Executable support--reserve recovery semantics.} Candidate prefixes are paired with actuator-realizable recovery continuations. A non-compensatory support--reserve admission uses one signed semantics across Safe, Near-Contact, and Contact: establish shared recovery support, preserve positive reserve, and repay negative recovery debt.
    \item \textbf{Candidate-conditioned constraint-native identifiability.} We show that candidate-invariant scene orientation is insufficient and formulate the current constraint-native audit in terms of signed candidate-induced active-constraint response, dimension-matched active-vs-nearest controls, and candidate-identity shuffling. The audit is offline and fail-closed, so mechanism evidence is not conflated with runtime capacity fitting.
    \item \textbf{Role isolation and unified evaluation.} RIFA separates absolute admission from relative ranking, while scene-disjoint Safe/Near-Contact/Contact strata evaluate nominal non-interference, critical-safety headroom, and post-contact recovery under the same policy.
\end{enumerate}

\section{Related Work}

\subsection{Post-Crash, Post-Impact, and Collision-Mitigation Planning}
Post-impact work demonstrates that safety does not end at the first collision. Robotaxi crash analyses highlight unsafe post-event evolution~\cite{koopman2024anatomy}, while post-impact vehicle-control studies address stabilization, controllability, and secondary-collision risk~\cite{wang2023integrated,parseh2023motion,wang2024post,yang2025post,ghosh2026post}. Collision-mitigation planners likewise consider unavoidable-impact severity and real-time evasive behavior~\cite{li2022vehicle,bosia2024real}. OC-RAP is complementary: it asks upstream whether a candidate action preserves a deployable recovery before contact, and it retains the same recovery semantics when contact has already occurred.

\subsection{Reachability, Viability, Barrier Functions, and Safety Filters}
Hamilton--Jacobi reachability, learned reachability, control barrier functions, and predictive/backup safety filters provide strong semantics for safe sets, invariance, reach-avoid behavior, and backup feasibility~\cite{bansal2017hamilton,bansal2021deepreach,lin2023generating,borquez2022parameter,ames2019control,koller2018learning,wabersich2021predictive,tearle2021predictive,kim2026backup,so2023solving}. Their existence guarantees do not automatically resolve our information-pattern question: a safe recovery may exist in every latent branch while no single post-prefix observation permits the deployed vehicle to select the branch-specific maneuver. OC-RAP therefore treats recovery choice itself as constrained by observation compatibility.

\subsection{Risk-Aware, Distributionally Robust, and Contingency Planning}
Risk-aware and distributionally robust planners use chance constraints, CVaR, ambiguity sets, scenario trees, or multimodal prediction~\cite{uryasev2000conditional,hakobyan2021wasserstein,safaoui2024distributionally,batkovic2020robust,nair2024predictive}. Contingency planners maintain multiple policies or fallback plans across futures~\cite{li2023marc,mustafa2024racp}, while formal driving models such as RSS impose structured safety rules~\cite{shalev2017formal}. OC-RAP shares the use of tail risk but differs in where the decision is constrained: the recovery option is fixed before aggregating latent futures that the deployed observation cannot distinguish.

\subsection{Learning-Based Planning, Benchmarks, and Reliable Admission}
Large-scale datasets, simulators, and closed-loop benchmarks have enabled data-driven autonomous planning~\cite{dosovitskiy2017carla,caesar2021nuplan,li2022metadrive,ettinger2021large,wilson2023argoverse,zhan2019interaction,althoff2017commonroad,krajewski2018highd,caesar2020nuscenes}. Recent surveys cover learning-based planning, risk assessment, and world models~\cite{ding2025understanding,lu2025risk,hu2025survey}. Calibration and conformal methods provide tools for statistically controlled admission under uncertainty~\cite{vovk2005algorithmic,stephen2021gentle,angelopoulos2024conformal,lekeufack2024conformal,strawn2023conformal}. In our formulation, calibration is downstream of a more basic question: which \emph{deployable recoverability property} should be admitted. We therefore first isolate the information pattern and the absolute source, then evaluate calibration and closed-loop performance.
\section{Method}
\label{sec:method}

OC-RAP separates four roles: executable recovery construction, observation-consistent recoverability evaluation, absolute admission, and relative candidate selection. Candidate-conditioned representation audits are offline and never receive privileged future information.

\subsection{Unified Recovery Planning}
At planning time $t$, let
\begin{equation}
 h_t=\{x^e_{t-H:t},\mathcal X_{t-H:t},\mathcal M,\Omega_t\}
\end{equation}
be the observable history. The upstream planner produces nominal prefix $a_0$ and candidate set $\cA_t=\{a_0,a_1,\ldots,a_N\}$. The deployed policy
\begin{equation}
 a_t^{\rm exec}=\pi_{\rm OC\mbox{-}RAP}(h_t,\cA_t)
\end{equation}
receives no Safe/Near/Contact identifier. Proposal size, recovery options, admission semantics, and fallback rules are shared across all regimes.

\subsection{Actuator-Realizable Recovery and OC-MERO}
Each candidate is paired with a finite recovery library $\cG=\{g_1,\ldots,g_L\}$. Desired acceleration/steering commands are projected into acceleration, jerk, steering-angle, and steering-rate envelopes before rollout; actuator feasibility is therefore enforced by construction rather than by a post-hoc veto.

For candidate $a$, the frozen recovery model provides latent roots $(z_k,p_k)$, post-prefix observation embeddings $e_k^o$, and root--option signed margins $\hat m_{k\ell}(a)$. Observation compatibility is
\begin{equation}
 C_{ij}(a)=\exp\!\left(-\frac{\|e_i^o-e_j^o\|_2^2/d_o}{\tau_o}\right),\qquad C_{ii}=1.
\end{equation}
OC-MERO first evaluates a \emph{fixed} recovery option over the compatible-root lower tail,
\begin{align}
 w_{ij}(a)&=\frac{C_{ij}(a)p_j(a)}{\sum_r C_{ir}(a)p_r(a)+\epsilon},\\
 q_{i\ell}(a)&=\LCVaR_{\beta}(\{\hat m_{j\ell}(a)\};\{w_{ij}(a)\}),
\end{align}
then permits each distinguishable observation anchor to select its best valid option,
\begin{align}
 r_i(a)&=\max_{\ell}q_{i\ell}(a),\\
 \widehat R_{\rm dep}(a)&=\LCVaR_{\alpha}(\{r_i(a)\};\{p_i(a)\}).
 \label{eq:rdep_main}
\end{align}
This ordering prevents hidden-root-dependent recovery choices from being counted as deployable capability.

\subsection{Support--Reserve Complementarity and RIFA}
Let $q_i^{\max}=\max_\ell q_{i\ell}$ and define the hard shared-support mass
\begin{equation}
 D_h(a)=\sum_i p_i\mathbf 1[q_i^{\max}(a)\ge0],
\end{equation}
with deployability coordinate $F(a)=\sigmoid(\widehat R_{\rm dep}(a))$. The absolute support--reserve predicate is
\begin{equation}
 \mathcal A_{\rm sr}(a\mid a_0)=
 \begin{cases}
 \mathbf 1[D_h(a)>0], & D_h(a_0)=0,\\
 \mathbf 1[D_h(a)>0]\mathbf 1[\widehat R_{\rm dep}(a)\ge0], & D_h(a_0)>0.
 \end{cases}
 \label{eq:srca_main}
\end{equation}
Thus positive deployability cannot compensate for destroying shared support. The same predicate realizes support establishment in low-headroom pre-contact states and reserve/debt semantics after support exists.

RIFA first forms the frozen relative proposal
\begin{equation}
 \cP_K=\TopK_{a\in\cA_t\setminus\{a_0\}}(\rho(a),K_p),\qquad K_p=5,
\end{equation}
applies Eq.~\eqref{eq:srca_main} only within $\cP_K$, and reranks only surviving candidates with the frozen relative evidence path. If none survives, OC-RAP executes $a_0$.

\subsection{Candidate-Conditioned Constraint-Native Orientation Audit}
\label{subsec:cnro_main}
A scene-wide orientation is insufficient when different candidate prefixes activate different interactions. The current V48.111 mechanism audit therefore keeps the first eight complete prefix states and uses only observation-side agent position, velocity, and footprint to form a constant-velocity continuation. For agent $i$ and prefix time $t_j$,
\begin{equation}
 h_{ij}^{a}=\|p_e^{a}(t_j)-p_i^{\rm CV}(t_j)\|-r_e-r_i.
\end{equation}
The candidate-induced signed constraint response is
\begin{equation}
 \Delta h_{ij}^{a}=h_{ij}^{a}-h_{ij}^{a_0}.
 \label{eq:delta_h_main}
\end{equation}
Positive $\Delta h$ means that the candidate moves the corresponding clearance constraint toward greater reserve; negative $\Delta h$ means increasing debt. The audit compares two equal-dimensional relational families: the two currently nearest agents versus the two candidate-specific minimum-clearance agents. Both append the same $2\times8\times2=32$ constraint-response coordinates to the same 156-D candidate pathway, producing a strict 188-D versus 188-D causal control. CNRO is audit-only in the current release; no audit score is inserted into runtime planning before preregistered Support and Reserve transfer gates are satisfied.

\begin{algorithm}[t]
\caption{OC-RAP runtime decision (CNRO remains offline)}
\label{alg:ocrap}
\begin{algorithmic}[1]
\STATE Generate nominal $a_0$ and candidates $\cA_t$ with the frozen upstream planner.
\STATE Realize each candidate with actuator-feasible recovery options.
\STATE Compute OC-MERO deployable recoverability $\widehat R_{\rm dep}(a)$ and shared-support mass $D_h(a)$.
\STATE Form frozen top-$K$ proposal $\cP_K$ from relative evidence.
\STATE Apply non-compensatory support--reserve admission in Eq.~\eqref{eq:srca_main} inside $\cP_K$.
\STATE Rerank surviving candidates with frozen relative evidence; execute $a_0$ if none survives.
\end{algorithmic}
\end{algorithm}

\section{Experiments}
\label{sec:experiments}

\subsection{Dataset and Three-Regime Protocol}
We use WOMD-derived scenes with Waymax-based closed-loop/counterfactual construction~\cite{ettinger2021large,gulino2023waymax}. Adaptation-train, development, certificate, and final-test roles are scene-disjoint. Safe, Near-Contact, and Contact are evaluation strata, not policy states. Safe measures nominal non-interference; Near-Contact emphasizes low-headroom pre-contact interaction and observation ambiguity; Contact contains contact/post-contact recovery, re-contact, and secondary-collision failure modes. No experiment in the V48.111 chain reconstructs or reselects the canonical dataset.

\begin{table}[t]
\centering
\caption{Primary closed-loop endpoints used for the submission run.}
\label{tab:regime_metrics}
\begin{tabular}{p{0.17\linewidth}p{0.75\linewidth}}
\toprule
Stratum & Primary endpoints \\
\midrule
Safe & collision/off-road scene rate; bounded nominal-utility preservation (NUP); intervention rate \\
Near-Contact & collision/off-road scene rate; 5th-percentile scene minimum clearance; 5th-percentile TTC; critical-TTC exposure duration; NUP; intervention rate \\
Contact & terminal clearance; normalized post-contact free-space AUC; clearance gain; escape rate; re-contact rate; secondary-overlap rate; stable-stop quality; overlap duration; off-road rate \\
\bottomrule
\end{tabular}
\end{table}

\subsection{Mechanism Evidence and Current Audit}
The mechanism-development chain is intentionally fail-closed. Earlier factor interventions identify support establishment and deployability improvement as the dominant complementary decision states. Subsequent representation audits show that a candidate-invariant scene-conditioned orientation can improve one population while rotating in another. V48.110 therefore tests candidate-specific active-agent topology; it provides local ordering signal but fails the preregistered full Support/Reserve transfer gate. The current V48.111 audit removes the principal capacity confound of that test by comparing 188-D active-pair and 188-D nearest-pair constraint-response features under the same strictly convex closed-form ridge solver.

The V48.111 null permutes candidate identity jointly with the candidate-induced topology/geometry within each scene-time group. This prevents the original candidate's active constraint from leaking into the shuffled control. Balanced and precision variants are retained for historical protocol compatibility, but identical raw-audit results are treated as duplicate protocol views rather than independent statistical replicates.

\subsection{Submission-Time Full Closed-Loop Evaluation}
The V48.111 repository is audit-oriented, so we additionally provide a submission evaluation harness that runs the frozen checkpoint and its frozen bucket-wise calibration end-to-end on Safe, Near-Contact, and Contact. CNRO remains offline: the harness does not insert its ridge solution into the deployed planner, does not finetune the checkpoint, and does not recalibrate thresholds. This distinction is essential for scientific validity: the final three-regime table evaluates the deployed frozen planner, while CNRO is reported as a representation/identifiability study.

\begin{table*}[t]
\centering
\caption{Submission closed-loop results. Fill directly from \texttt{V48.111-SUBMISSION-THREE-REGIME-SUMMARY.csv}; no metric should be copied from mechanism-audit tables.}
\label{tab:main_results}
\setlength{\tabcolsep}{4pt}
\begin{tabular}{llcccccccc}
\toprule
Variant & Regime & Collision$\downarrow$ & Off-road$\downarrow$ & NUP$\uparrow$ & Interv.$\downarrow$ & Clear.$\uparrow$ & TTC$\uparrow$ & Re-contact$\downarrow$ & Escape$\uparrow$ \\
\midrule
Balanced & Safe & -- & -- & -- & -- & -- & -- & -- & -- \\
Balanced & Near & -- & -- & -- & -- & -- & -- & -- & -- \\
Balanced & Contact & -- & -- & -- & -- & -- & -- & -- & -- \\
Precision & Safe & -- & -- & -- & -- & -- & -- & -- & -- \\
Precision & Near & -- & -- & -- & -- & -- & -- & -- & -- \\
Precision & Contact & -- & -- & -- & -- & -- & -- & -- & -- \\
\bottomrule
\end{tabular}
\end{table*}

\subsection{Reporting Discipline}
We distinguish three evidence levels: (i) information preservation, (ii) finite-sample accessibility under a fixed probe, and (iii) population-shared deployable action ordering. A lower audit loss is not a substitute for the third criterion. Likewise, final closed-loop results are not used to retroactively redefine the audit gates. External baselines are meaningful only after the deployed source/policy is frozen for the submitted model.

\section{Conclusion}
\label{sec:conclusion}
OC-RAP treats recovery as a deployable action property under partial observation rather than a branch-wise oracle label. OC-MERO constrains recovery choice by observation compatibility, actuator projection makes recovery continuations executable, support--reserve complementarity provides one signed semantics before and after contact, and RIFA isolates absolute admission from relative ranking. The current representation evidence further sharpens the invariance principle: recovery orientation should be candidate-conditioned and constraint-native because the candidate itself changes which interaction limits recoverability. V48.111 tests this claim through nominal-zero signed constraint response with a dimension-matched active-vs-nearest causal control while keeping the deployed planner frozen. This separation between runtime mechanism and offline identifiability is deliberate. It allows Safe, Near-Contact, and Contact to remain one policy and one recoverability object---positive reserve before boundary loss and negative debt after violation---without regime routing or privileged future information.

\appendices
\section{Detailed Formulation and Reproducibility}
\label{sec:appendix}

\subsection{Executable Recovery Projection}
A recovery option outputs desired acceleration and steering $(a_s^{\rm des},\delta_s^{\rm des})$. With step $\Delta t$, acceleration bounds $[a_{\min},a_{\max}]$, jerk limit $j_{\max}$, steering bound $\delta_{\max}$, and steering-rate bound $\dot\delta_{\max}$,
\begin{align}
 a_s &= \Pi_{[\max(a_{\min},a_{s-1}-j_{\max}\Delta t),\,\min(a_{\max},a_{s-1}+j_{\max}\Delta t)]}(a_s^{\rm des}),\\
 \delta_s &= \Pi_{[\max(-\delta_{\max},\delta_{s-1}-\dot\delta_{\max}\Delta t),\,\min(\delta_{\max},\delta_{s-1}+\dot\delta_{\max}\Delta t)]}(\delta_s^{\rm des}).
\end{align}
This retains option identity while enforcing actuator feasibility by construction. It avoids interpreting a controllable command discontinuity as an environmental feasibility failure.

\subsection{OC-MERO Lower-Tail Operator}
For normalized weights $w_i$, the implementation uses stable sorting and fractional lower-tail mass. An equivalent variational form is
\begin{equation}
 \LCVaR_{\alpha}(\{s_i\};\{w_i\})=\max_{\nu\in\mathbb R}\left[\nu-\frac{1}{\alpha}\sum_iw_i[\nu-s_i]_+\right].
\end{equation}
For a fixed observation anchor $i$, OC-MERO first aggregates each option over observation-compatible roots using $\beta$-LCVaR, then takes the best valid option, and finally aggregates observation anchors using $\alpha$-LCVaR. The current mechanism chain uses $\alpha=\beta=0.2$ and fixed top-$M=8$ compatibility sparsification where enabled.

A hidden-root oracle instead computes
\begin{equation}
 \widehat R_{\rm orc}(a)=\LCVaR_{\beta}(\{\max_\ell \hat m_{k\ell}(a)\};\{p_k\}),
\end{equation}
and the oracle gap is $\widehat G(a)=\widehat R_{\rm orc}(a)-\widehat R_{\rm dep}(a)$.

\subsection{Native Recovery Certificate and Complementarity}
Let $q_i^{\max}=\max_\ell q_{i\ell}$. Besides $D_h$ and $F$, the implementation records
\begin{align}
 D_s(a)&=\sum_i p_i\sigmoid(q_i^{\max}(a)/\tau_q),\\
 Q_G(a)&=\exp(-[\widehat G(a)]_+).
\end{align}
The native certificate is $c(a)=[D_h,F,D_s,Q_G]$. $D_h=0$ is the shared-support existence boundary and $F=0.5$ is the $\widehat R_{\rm dep}=0$ boundary. The positive absolute predicate is non-compensatory: a candidate that destroys shared support cannot be rescued by another coordinate.

\subsection{Teacher Physical Margin and Structural Contract}
The pre-structural teacher physical margin is a signed minimum over active normalized constraints,
\begin{equation}
 m_{\rm phys}^{\star}(a,z,g)=\min_{q\in\mathcal Q_{\rm active}}\frac{b_q(\xi_{a,z,g})}{s_q}.
\end{equation}
The stored target is not unconditionally this pure minimum. The current generator may apply ordered structural operations: non-hidden \texttt{post\_contact\_stabilize}, \texttt{yield\_rejoin}, and \texttt{pull\_over} modes may receive a floor at $0.6$; blocked \texttt{yield\_rejoin} may be capped at $-0.8$; active \texttt{avoid\_secondary} may receive a floor at $0.9$; and deliberately mined hidden branches may use branch-specific hard replacement. Structural metadata is teacher/audit-side only. We therefore distinguish physical reserve/debt semantics from structural deployability semantics and do not describe the authoritative stored target as a pure physical slack minimum.

\subsection{V48.111 Constraint-Native Recovery Orientation Audit}
\label{app:cnro}
The V48.111 audit is designed to isolate one hypothesis: \emph{the missing invariant is candidate-conditioned signed constraint response rather than raw active-agent identity}. It changes no deployed planner parameter.

\subsubsection{Observation-only active geometry}
The historical candidate pathway contains an 80-D prefix-state block. V48.111 uses the first eight complete 9-D prefix states and excludes the truncated ninth record, producing the registered early-action window. For each currently observed agent, current relative position, velocity, length, and width define a constant-velocity continuation. With circle radii $r_e,r_i$,
\begin{equation}
 h_{ij}^{a}=\|p_e^{a}(t_j)-p_i^{\rm CV}(t_j)\|-r_e-r_i.
\end{equation}
For each agent,
\begin{equation}
 h_i^{a}=\min_j h_{ij}^{a}.
\end{equation}
The active pair is the two agents with smallest $h_i^a$; the control pair is the two currently nearest agents. Neither selector uses teacher future, regime identity, or learned routing.

\subsubsection{Nominal-zero constraint response}
For candidate $a$ and nominal $a_0$,
\begin{equation}
 \Delta h_{ij}^{a}=h_{ij}^{a}-h_{ij}^{a_0}.
\end{equation}
To preserve boundary context without moving the physical zero boundary, V48.111 uses two channels per selected agent/time:
\begin{equation}
 \phi_{ij}^{(1)}=\frac{\Delta h_{ij}^{a}}{s_h},\qquad
 \phi_{ij}^{(2)}=\frac{\Delta h_{ij}^{a}h_{ij}^{a_0}}{s_h^2},
\end{equation}
where $s_h$ is a train-only family-independent RMS scale. With two agents and eight times, each relational family adds 32 coordinates to the same 156-D raw candidate pathway, so
\begin{equation}
 d_{\rm nearest}=d_{\rm active}=156+2\times8\times2=188.
\end{equation}
This removes the V48.110 1716-D versus 3588-D capacity mismatch from the active-vs-nearest contrast.

\subsubsection{Closed-form orientation probe}
Support and Reserve are audited separately with the same float64 class-balanced ridge objective,
\begin{equation}
 \hat w=\argmin_w \sum_n \omega_n(y_n-x_n^\top w)^2+\lambda\|w\|_2^2,
\end{equation}
with $\lambda=1/N$. The problem is strictly convex and solved in closed form; there is no epoch, learning-rate, patience, rank, or width sweep. Consequently, optimization failure cannot explain a negative result.

\subsubsection{Candidate-identity shuffle}
Because the selected active pair is candidate-dependent, the null permutes the entire candidate identity inside each scene-time group rather than shuffling the 156-D response alone. Candidate response, active selection, nearest selection, and derived geometry move together. Scene context remains fixed. This tests whether candidate identity plus candidate-induced constraint response carries action information without leaking the original candidate's topology into the null.

\subsubsection{Preregistered interpretation}
Support/Reserve use the historical absolute AUC, true-minus-shuffle, and top-1 gates. In addition, V48.111 requires an active-switch increment: active-pair must outperform the equal-dimensional nearest-pair control across the registered Near/Contact development/certificate roles. Only joint Support+Reserve transfer and matched active-minus-nearest improvement can authorize promotion of a candidate-conditioned constraint-native carrier. Otherwise the fixed-CV circle geometry family is closed rather than widened or tuned.

\subsection{Counterfactual Equivalence-Partition Transport}
Let future instances be partitioned into equivalence classes $e$ from branch-defining semantic recipe and, for stochastic/augmented branches, stored exogenous realization. Let $c_{ke}$ and $n_{je}$ be candidate and nominal probability mass assigned to roots $k$ and $j$. For a shared class,
\begin{equation}
 \Pi_e(k,j)=\min(c_e,n_e)\frac{c_{ke}}{c_e}\frac{n_{je}}{n_e},\qquad \Pi=\sum_e\Pi_e.
\end{equation}
This transport is invariant to root-slot permutation and represents split/merge correspondence. Missing exogenous fingerprints are unresolved rather than silently matched. It is used only for offline identifiability diagnostics.

\subsection{Closed or Diagnostic-Only Families}
The mechanism record provides negative evidence that should not be converted into extra Main components. Closed/frozen primary fixes include option-wise scalar/typed gains, class-local/path-stop branches, repeated CV/CA and ball/box/hull/anchor/jerk sweeps, high-order signed-viability sweeps, repeated truth-floor/cap inverse refinement, frozen structured-tail width/rank/depth sweeps, generic root/action MLP replacement, proposal/top-$K$ expansion, premature boundary transport, relative-ranker modification before source freeze, regime routers/experts/thresholds/budgets, privileged-future inputs, and dataset reconstruction. V48.110 additionally closes raw active-agent-token topology as a complete mechanism; V48.111 tests only the constraint-native response hypothesis under an equal-dimensional causal control.

\subsection{Code-Aligned Closed-Loop Metrics}
\begingroup\footnotesize\sloppy
Safe reports \path{collision_scene_rate}, \path{offroad_scene_rate}, \path{closed_loop_bounded_NUP}, and \path{intervention_rate}. Near-Contact additionally reports \path{scene_min_clearance_m_p05}, \path{scene_ttc_s_p05}, and \path{critical_ttc_exposure_duration_s}. Contact reports \path{post_contact_terminal_clearance_m}, \path{post_contact_free_space_auc_normalized_m}, \path{post_contact_clearance_gain_m}, \path{post_contact_escape_scene_rate}, \path{recontact_scene_rate}, \path{secondary_overlap_scene_rate}, \path{new_stable_stop_quality_scene_rate}, and \path{post_contact_overlap_duration_s}. The submission harness writes these values to a single machine-readable summary rather than mixing mechanism probes with final closed-loop metrics.
\endgroup

\bibliographystyle{IEEEtran}
\bibliography{post-collision}

\end{document}
